\section{Súčasný stav riešenej problematiky}

Na trhu existuje množstvo platforiem zameraných na organizáciu a objavovanie podujatí. Dominantným hráčom v tejto oblasti je služba \textbf{Facebook Events}, vyvíjaná spoločnosťou Meta. Vďaka zázemiu veľkej technologickej spoločnosti ide o komplexný systém, ktorý pokrýva široké spektrum potrieb používateľov.

\subsection{Facebook Events}
Facebook Events umožňuje používateľom vytvárať vyhradené stránky pre konkrétne príležitosti a zasielať verejné alebo súkromné 
pozvánky.
Pre online firmy predstavujú tieto podujatia efektívny nástroj
na propagáciu predajov, osláv alebo dôležitých míľnikov. 
Systém zahŕňa históriu podujatí, možnosť označovania obľúbených aktivít a pokročilé filtrovanie.

\begin{itemize}[noitemsep]
    \item \textbf{Výhody:} Obrovská používateľská základňa, vysoká miera integrácie so sociálnou sieťou, 
    prepracované algoritmy na predikciu záujmov a množstvo funkcií dostupných bez priamych poplatkov za vytvorenie podujatia.
    \item \textbf{Nevýhody a obmedzenia:} 
    \begin{itemize}[noitemsep]
        \item \textbf{Informačné zahltenie a notifikačná únava:} Rozhranie je preplnené reklamami. Nadmerné množstvo systémových upozornení
         často vedie k tomu, že používatelia dôležité pozvánky na podujatia prehliadnu alebo ignorujú.
        \item \textbf{Nízka miera kontroly nad dosahom:} Organizátor nemá záruku, že sa informácia o podujatí dostane ku všetkým sledovateľom.
         Algoritmus uprednostňuje platenú reklamu a interaktívny obsah, čo znevýhodňuje menšie lokálne a komunitné aktivity.
        \item \textbf{Uzavretosť ekosystému:} Plné využitie funkcií vyžaduje existenciu Facebook profilu, čo odrádza časť modernej populácie.
        \item \textbf{Obmedzená vizualizácia:} Mapa podujatí nie je primárnym ovládacím prvkom a je dostupná až po manuálnom rozkliknutí, 
        čo znižuje intuitívnosť pri hľadaní aktivít v okolí.
        \item \textbf{Absencia priamej komunikácie:} Diskusia v rámci podujatia prebieha formou príspevkov na nástenke,
         chýba však funkcia reálneho skupinového chatu pre účastníkov v reálnom čase.
    \end{itemize}
\end{itemize} \cite{facebook_events, bigcommerce_facebook_events , stewart2020facebook}

\subsection{Meetup}
Platforma Meetup je celosvetovo známa platforma sa primárne zameriava na budovanie dlhodobých komunít a záujmových skupín. 
Na rozdiel od Facebooku, kde sú podujatia často jednorazové, Meetup motivuje používateľov k pravidelnému stretávaniu sa na základe spoločných koníčkov, 
ako sú technológie, šport či jazykové vzdelávanie.
V porovnaní s Facebook Events ponúka Meetup detailnejšie komunitné funkcie, avšak väčšina z nich je podmienená plateným členstvom:

\begin{itemize}[noitemsep]
    \item \textbf{Rozšírená interakcia:} Platforma umožňuje prehliadať zoznamy účastníkov, pristupovať k ich profilom a využívať pokročilé 
    systémy komentárov a správ, čo uľahčuje nadväzovanie kontaktov ešte pred samotným podujatím.
    \item \textbf{Personalizácia} Algoritmus navrhuje podujatia a skupiny na základe histórie aktivít a preferencií používateľa,
     pričom prioritizuje lokálne komunity.
    \item \textbf{Platený model (Monetizácia):} Najväčšou nevýhodou platformy je, že kľúčové funkcie pre organizátorov 
    (vytváranie podujatí, správa členov) aj pokročilé funkcie pre účastníkov sú spoplatnené. V rámci prémiového balíka získavajú používatelia:
    \begin{itemize}[noitemsep]
        \item prístup k úplným zoznamom a profilom účastníkov,
        \item prioritný prístup k podujatiam a čakacím listinám,
        \item neobmedzené zasielanie správ ostatným členom,
        \item rozhranie bez reklám a včasné upozornenia na nové udalosti.
    \end{itemize}
    \item \textbf{Problém s účasťou:} Podobne ako pri iných bezplatných registráciách, aj Meetup čelí vysokej miere tzv. \textit{no-shows} (používatelia, ktorí potvrdia účasť, ale neprirodzene absentujú), čo sťažuje plánovanie kapacity pre organizátorov.
\end{itemize}

\noindent \textbf{Kritické zhodnotenie:} Hoci Meetup poskytuje kvalitnejšie nástroje na sieťovanie než Facebook, jeho uzavretosť za platobnú bránu a slabšia SEO viditeľnosť pre nečlenov otvárajú priestor pre vývoj prístupnejších, bezplatných alternatív \cite{somerfleck2021meetup, meetup_agile_bratislava}.
\begin{itemize}[noitemsep]
    \item \textbf{Problém „No-shows“ (Neprihlásení účastníci):} Štatistiky z praxe ukazujú, že v priemere až 50 \% 
    používateľov, ktorí potvrdia účasť, na udalosť nakoniec nepríde. 
    \item \textbf{Vysoké prevádzkové náklady:} Pre organizátorov predstavuje Meetup finančnú záťaž (poplatky od cca 25 \$ mesačne), 
    pričom základné balíky výrazne obmedzujú možnosti brandingu, vkladanie multimédií (video, vlastné logá) či prispôsobenie dizajnu skupiny.
    \item \textbf{Geografické obmedzenia vyhľadávania:} Systém je primárne postavený na vyhľadávaní podľa PSČ alebo mesta.
     To drasticky limituje dosah udalostí, ktoré sú organizované online alebo cília na globálne publikum, 
     nakoľko používatelia nevidia relevantné výsledky mimo svojho bezprostredného okolia.
    \item \textbf{Pasivita a „Ghost groups“:} Na platforme sa nachádza veľké množstvo neaktívnych skupín, ktoré slúžia len na 
    zber e-mailových adries bez reálnej snahy organizovať stretnutia.
    \item \textbf{Kvalita a bezpečnosť lokalít:} Mnohé stretnutia sa odohrávajú v baroch alebo reštauráciách, čo prináša problémy s hlukom, 
    nevhodným osvetlením alebo dokonca neprofesionálnym správaním účastníkov (alkohol). 
    Verejné inštitúcie (knižnice) sú často pre neformálne skupiny nedostupné.
    \item \textbf{Chýbajúca automatizácia časových pásiem:} Meetup nedostatočne rieši automatický prepočet časových pásiem pri medzinárodných online podujatiach, čo núti organizátorov k manuálnym opravám.
\end{itemize}


\subsection{AllEvents.in}
\textit{AllEvents} sa prezentuje ako globálna platforma na objavovanie udalostí, ktorá pokrýva viac než 40 000 miest po celom svete \cite{allevents_web}. Platforma sa zameriava na agregáciu širokého spektra podujatí – od komunitných stretnutí až po veľké komerčné koncerty a festivaly.

\begin{itemize}[noitemsep]
    \item \textbf{Deklarované funkcie:} Personalizované odporúčania na základe záujmov, sledovanie aktivít priateľov, bezpapierová registrácia (\textit{check-in}) a priama rezervácia vstupeniek v rámci aplikácie.
    \item \textbf{Kategorizácia a distribúcia:} Platforma vyniká prehľadným členením do kategórií, možnosťou exportu do kalendára a funkciou priameho kontaktovania organizátora (\textit{Message to host}). 
    \item \textbf{Identifikované nedostatky a technické bariéry:} 
    \begin{itemize}[noitemsep]
        \item \textbf{Kritická nestabilita aplikácie:} Pri testovaní mobilnej aplikácie boli zaznamenané závažné chyby v procese autentifikácie (\textit{Login failed}) a zlyhania pri verifikácii používateľa cez mobilné číslo, čo znemožňuje plnohodnotné využívanie služieb \cite{allevents_play}.
        \item \textbf{Lokalizačné problémy:} Napriek deklarovanému globálnemu dosahu systém v niektorých prípadoch nedokáže korektne identifikovať menšie mestá, čo núti používateľa k prehliadaniu v režime hosťa (\textit{Guest mode}).
        \item \textbf{Absencia sociálnej interakcie:} Hoci platforma umožňuje nákup lístkov a sledovanie udalostí, postráda funkciu reálneho skupinového chatu medzi účastníkmi a transparentné zobrazenie zoznamu prihlásených osôb.
        \item \textbf{Zameranie obsahu:} Databáza je orientovaná primárne na veľké komerčné a spoločenské akcie, čím opomína drobné, spontánne komunitné aktivity.
    \end{itemize}
\end{itemize}



\subsection{Hapsy}
Aplikácia \textit{Hapsy} sa profiluje ako agregátor podujatí , ktorého cieľom je zjednotiť ponuku z rôznych platforiem (Facebook, Eventbrite, Meetup, GoOut) do jedného interaktívneho zobrazenia na mape. Jej hlavnou ambíciou je eliminovať potrebu prepínania medzi viacerými aplikáciami.

\begin{itemize}[noitemsep]
    \item \textbf{Kľúčové funkcie:} Agregácia dát z externých zdrojov, personalizované upozornenia na základe preferencií používateľa, ukladanie obľúbených udalostí a prehľadná mapa mesta.
    \item \textbf{Identifikované obmedzenia a používateľská kritika:} 
    \begin{itemize}[noitemsep]
        \item \textbf{Geografická limitácia:} Podľa spätnej väzby používateľov trpí aplikácia nedostatkom 
        lokálneho obsahu mimo hlavných miest. Algoritmus sa zameriava na veľké komerčné databázy, čím opomína komunitný život v menších regiónoch.
        \item \textbf{Absencia tvorby obsahu:} Systém funguje primárne ako prehliadač.
         Bežný používateľ nemá možnosť vytvoriť vlastnú udalosť, čo znemožňuje organizáciu spontánnych aktivít.
        \item \textbf{Technická nedokonalosť:} Používateľské recenzie naznačujú prítomnosť kritických chýb v rozhraní.
        \item \textbf{Obmedzená interakcia:} Podobne ako pri iných agregátoroch, aj tu chýba sociálna vrstva - možnosť priamej komunikácie medzi účastníkmi alebo správa tímov.
    \end{itemize}
\end{itemize}\cite{hapsy_play}



\subsection{Mapibara}
Platforma \textit{Mapibara} predstavuje minimalistický prístup k objavovaniu podujatí,
pričom kladie absolútny dôraz na mapové rozhranie a anonymitu používateľa .
Na rozdiel od komplexných sociálnych sietí sa zbavuje povinných registrácií a zameriava sa na okamžitý prístup k informáciám.

\begin{itemize}[noitemsep]
    \item \textbf{Silné stránky dizajnu:} Extrémne jednoduché a prehľadné používateľské rozhranie, kde mapa zaberá celú plochu obrazovky. 
    Absencia prihlasovacieho procesu odstraňuje vstupnú bariéru pre nových používateľov.
    \item \textbf{Identifikované kritické nedostatky:} 
    \begin{itemize}
        \item \textbf{Absencia obsahu:} Pri praktickom testovaní platforma vykazuje takmer nulovú naplnenosť podujatiami, 
        čo ju robí v súčasnom stave nepoužiteľnou.
        \item \textbf{Chýbajúca interakcia:} Kvôli absencii používateľských profilov nie je možné s účastníkmi 
        komunikovať, spravovať kapacity ani riešiť logistiku (napr. spomínané rozdelenie do tímov pri športe).
        \item \textbf{Pasivita systému:} Platforma slúži len ako statický prehliadač bodov na mape bez možnosti aktívneho zapojenia komunity.
    \end{itemize}
\end{itemize}\cite{mapibara}

\noindent \textbf{Ponaučenie pre vlastný návrh:} Analýza Mapibary ukazuje, že samotná vizualizácia na mape nestačí. Pre úspech aplikácie je nevyhnutné skombinovať \textbf{jednoduchosť mapy} so \textbf{sociálnou interakciou} (prihlásenie, chat), aby používatelia mali dôvod podujatia nielen hľadať, ale aj vytvárať a navštevovať.

\subsection{Zhrnutie analýzy súčasného stavu}

Na základe podrobnej analýzy dominantných platforiem (Facebook Events, Meetup, AllEvents, Hapsy a Mapibara) je možné definovať súbor štandardných funkcií, ktoré tvoria základ moderných systémov na správu podujatí, ale zároveň identifikovať kritické funkčné medzery.

\paragraph{Štandardné funkcie prítomné v analyzovaných riešeniach:}
Väčšina spomenutých aplikácií disponuje nasledujúcim funkčným základom:
\begin{itemize}[noitemsep]
    \item \textbf{Objavovanie a vizualizácia:} Interaktívna mapa s podujatiami, vyhľadávanie podľa kategórií a pokročilé filtrovanie.
    \item \textbf{Informačná vrstva:} Detailný popis podujatia, fotogaléria, profil organizátora a zdieľanie  na iné platformy.
    \item \textbf{Personalizácia:} História navštívených udalostí, systém odporúčaných podujatí a push notifikácie.
    \item \textbf{Interakcia:} Možnosť vytvorenia vlastnej verejnej alebo súkromnej udalosti po predchádzajúcej registrácii.
\end{itemize}

\paragraph{Identifikácia funkčných medzier (Inovačný potenciál):}
Napriek širokej škále funkcií analýza odhalila oblasti, ktoré súčasné riešenia pokrývajú nedostatočne, 
čo vytvára priestor pre vývoj nového, špecializovaného riešenia:

\begin{itemize}[noitemsep]
    \item \textbf{Dynamický manažment kapacity:} Absencia automatizovaného systému náhradníkov. Väčšina platforiem neponúka riešenie pre prípady, kedy je kapacita naplnená a nie je možné automaticky osloviť ďalšieho záujemcu v poradí pri uvoľnení miesta.
    \item \textbf{Verifikácia reálnej účasti:} Kritický nedostatok nástrojov na potláčanie efektu (používatelia, ktorí potvrdia účasť, ale neprídu). 
    \item \textbf{Komunitná logistika:} Absentujú nástroje na správu spoločného vybavenia  a funkcia reálneho skupinového chatu pre všetkých účastníkov v reálnom čase.
    \item \textbf{Moderovanie a bezpečnosť:} Pri raste používateľmi generovaného obsahu chýbajú v mnohých aplikáciách citlivé 
    nástroje na nahlasovanie a včasné blokovanie nevhodného obsahu alebo používateľov.
    \item \textbf{Regionálna relevantnosť:} Nízka hustota obsahu v menších mestách a regiónoch mimo hlavných metropol, 
    čo je dôsledkom orientácie veľkých platforiem na komerčné a masové podujatia.
\end{itemize}

\subsection{Syntéza poznatkov a východiská pre vlastný návrh}
Na základe vykonanej analýzy existujúcich riešení je možné konštatovať, že hoci trh ponúka robustné platformy, 
žiadna z nich nepokrýva komplexne potreby malých, spontánnych komunít s dôrazom na logistiku a potláčanie negativity.
Pre lepšiu prehľadnosť sú kľúčové rozdiely medzi analyzovanými platformami a navrhovaným riešením zhrnuté v Tabuľke \ref{tab:comparison_market}.

\begin{table}[h!]
\centering
\caption{Porovnanie kľúčových funkcií existujúcich platforiem a navrhovanej aplikácie}
\label{tab:comparison_market}
\begin{tabular}{|l|c|c|c|c|c|c|}
\hline
\textbf{Funkcia / Platforma} & \textbf{FB} & \textbf{Meetup} & \textbf{AllEvents} & \textbf{Hapsy} & \textbf{Mapibara} \\ \hline
Mapa ako primárne UI & Nie & Čiastočne & Áno & Áno & Áno \\ \hline
Tvorba udalostí & Áno & Platené & Áno & Nie & Nie \\ \hline
Skupinový chat & Nie & Áno & Nie & Nie & Nie \\ \hline
Manažment náhradníkov & Nie & Áno & Nie & Nie & Nie \\ \hline
Ochrana „No-show“ & Nie & Nie & Nie & Nie & Nie \\ \hline
Zameranie na komunity & Áno & Áno & Nie & Nie & Nie  \\ \hline
\end{tabular}
\end{table}


\paragraph{Definovanie smerovania vlastného riešenia}
Zistené nedostatky konkurenčných riešení priamo formujú funkčné požiadavky na moju aplikáciu. 
Vlastné riešenie sa bude odlišovať predovšetkým v týchto pilieroch:

\begin{enumerate}[noitemsep]
    \item \textbf{Interaktivita a operatívna komunikácia:} Integrácia komunikačného modulu priamo do detailu udalosti, čo umožňuje riešiť zmeny v reálnom čase a zvyšuje priamu angažovanosť účastníkov.
    \item \textbf{Zvýšenie kredibility a záväznosti:} Implementácia overovacích mechanizmov (napr. potvrdzovanie prítomnosti cez polohu alebo systém potvrdzovania krátko pred akciou), ktoré motivujú k reálnej účasti.
    \item \textbf{Minimalizmus a organické objavovanie:} Transparentné mapové rozhranie bez hraníc miest a zbytočného balastu, ktoré umožňuje používateľom objavovať udalosti prirodzeným scrollovaním.
    \item \textbf{Dostupnosť a odstránenie bariér:} Eliminácia finančných nákladov pre organizátorov a zabezpečenie okamžitého prehľadu o dianí v okolí pre všetkých používateľov.
    \item \textbf{Natívna podpora online podujatí:} Systémové riešenie pre digitálne udalosti s automatickou prácou s časom a optimalizáciou pre vyhľadávače.
\end{enumerate}

Tento prístup umožní vytvoriť nástroj, ktorý nie je len „nástenkou podujatí“, ale živým ekosystémom pre lokálne komunity.

\noindent Zhrnutím možno konštatovať, že existuje dopyt po riešení,
ktoré by premostilo priepasť medzi statickou informáciou o udalosti a aktívnou logistickou podporou komunitných aktivít.













